% © 2026 Ing. David Jaroš — CC BY-NC-ND 4.0
%
% This work is licensed under a Creative Commons Attribution-NonCommercial-NoDerivatives
% 4.0 International License (CC BY-NC-ND 4.0).
%
% License History: Earlier drafts (up to v0.3) were released under CC BY 4.0.
% From v0.4 onward, all material is released under CC BY-NC-ND 4.0 to protect
% the integrity of the theoretical work during ongoing academic development.
%
% See LICENSE.md for full license text.

% UBT-AI-PROVENANCE-BEGIN
% schema: ubt-ai-provenance/v1
% tier: B_machine_verified
% ai_assistance: disclosed
% human_review: machine-verification
% editorial_responsibility: Ing. David Jaroš
% policy: ../../AI_PROVENANCE.md
% notice: Machine-verified against named sources or verifiers; individual attestation is not claimed.
% UBT-AI-PROVENANCE-END

% -----------------------------------------------------------------------
% STATUS: HISTORICAL / COMPLEMENTARY — NOT THE CURRENT CANONICAL SU(3) DERIVATION
%
% This file summarises the involution-based approach to SU(3) colour (Theorems
% G.A–G.D).  It is preserved for provenance and priority.
%
% Key limitation (do not modify this notice):
%   The involution construction PROVED the carrier selection
%   Vc ≅ C^3 (a genuine algebraic result), but the unitary/Yang–Mills
%   dynamics on that carrier were POSTULATED, not derived from the UBT action.
%   GAP-SU3-DYN therefore remains OPEN.
%
% The current canonical derivation derives the Hermitian form and volume form
% algebraically from the biquaternionic adjoint and quaternion multiplication:
%   canonical/su3_derivation/su3_stabilizer_exterior_fock.tex
% -----------------------------------------------------------------------

\ifdefined\INCLUDEMODE
\else
  \documentclass[12pt,a4paper]{article}
  \usepackage[a4paper, margin=2.5cm]{geometry}
  \usepackage{amsmath, amssymb, amsthm}
  \usepackage{mathtools}
  \usepackage{hyperref}
  \usepackage{xcolor}
  \usepackage{mdframed}
  \usepackage{booktabs}

  \newtheorem{theorem}{Theorem}[section]
  \newtheorem{lemma}[theorem]{Lemma}
  \newtheorem{proposition}[theorem]{Proposition}
  \newtheorem{corollary}[theorem]{Corollary}
  \theoremstyle{definition}
  \newtheorem{definition}[theorem]{Definition}
  \theoremstyle{remark}
  \newtheorem{remark}[theorem]{Remark}

  \newmdenv[backgroundcolor=yellow!20, linecolor=orange, linewidth=1.5pt]{gapbox}
  \newmdenv[backgroundcolor=green!15, linecolor=green!60!black, linewidth=1.5pt]{resultbox}
  \newmdenv[backgroundcolor=blue!8, linecolor=blue!50, linewidth=1.5pt]{insightbox}

  \title{\textbf{SU(3) from Involutions on $\mathbb{C}\otimes\mathbb{H}$}\\
         \large Summary of Theorems G.A--G.D from Appendix~G}
  \author{David Jaro\v{s}}
  \date{2026-03-06}

  \begin{document}
  \maketitle

  \begin{abstract}
  We provide a clean summary of the derivation of $SU(3)_{\mathrm{color}}$
  from the biquaternion algebra $\mathbb{C}\otimes\mathbb{H}$ via involutions.
  This document collects Theorems G.A--G.D from
  \texttt{appendix\_G\_internal\_color\_symmetry.tex} in a compact form
  suitable for the main paper.  The Lie algebra $\mathfrak{su}(3)$, the
  fundamental representation (quarks), the adjoint representation (gluons),
  and the separation from electroweak are all proved.
  The sole remaining gap --- color confinement --- is a Clay Millennium
  Prize Problem; qualitative mechanism exists but quantitative derivation
  is future work.
  \end{abstract}
\fi

%──────────────────────────────────────────────────────────────────────────────
\section{Overview and Prior Approaches}
\label{sec:overview}
%──────────────────────────────────────────────────────────────────────────────

Earlier attempts to derive $SU(3)_c$ in UBT used the three imaginary
quaternionic axes $\{I, J, K\}$ directly.  This approach was correct
in spirit but fell short: the generators $\{I, J, K\}$ span a space
isomorphic to $\mathfrak{su}(2)$, not $\mathfrak{su}(3)$
($SU(3)$ has rank 2 and 8 generators, while $\{I,J,K\}$ gives rank 1
and 3 generators).

The correct derivation, given in \texttt{appendix\_G\_internal\_color\_symmetry.tex},
uses the \emph{involution approach}:

\begin{insightbox}
\textbf{Key insight:}
$\mathbb{C}\otimes\mathbb{H}$ has $\dim_\mathbb{R} = 8$.
Defining the carrier space $V_c = \ker(P_2 + 1)$ (the eigenspace of the
quaternionic conjugation involution with eigenvalue $-1$), we get a
3-dimensional complex space.  The automorphisms of $V_c$ preserving the
UBT inner product form $SU(3)$.
\end{insightbox}

%──────────────────────────────────────────────────────────────────────────────
\section{The Involution Construction}
\label{sec:involution}
%──────────────────────────────────────────────────────────────────────────────

\begin{definition}[Involutions on $\mathbb{C}\otimes\mathbb{H}$]
Let $\mathcal{B} = \mathbb{C}\otimes\mathbb{H}$ with basis
$\{1, I, J, K, i, iI, iJ, iK\}$ (where $i^2 = -1$, $I^2 = J^2 = K^2 = -1$,
$IJ = K$, etc.):
\begin{align}
  P_1 &:\quad i \mapsto -i
  \quad (\text{complex conjugation}), \\
  P_2 &:\quad I,J,K \mapsto -I,-J,-K
  \quad (\text{quaternionic conjugation}).
\end{align}
\end{definition}

\begin{definition}[Triplet carrier space]
\label{def:triplet}
The $(-1)$-eigenspace of $P_2$ on $\mathcal{B}$:
\begin{equation}
  V_c \;=\; \ker(P_2 + 1)
       \;=\; \mathrm{span}_\mathbb{C}\{I, J, K\}
       \;\cong\; \mathbb{C}^3.
\end{equation}
This is a 3-dimensional complex vector space --- the quark triplet.
\end{definition}

%──────────────────────────────────────────────────────────────────────────────
\section{Main Theorems}
\label{sec:theorems}
%──────────────────────────────────────────────────────────────────────────────

\begin{theorem}[G.A: Lie Algebra $\mathfrak{su}(3)$ --- PROVED]
\label{thm:G_A}
The group of automorphisms of $V_c$ preserving the UBT inner product
$\langle u, v\rangle = \mathrm{Sc}[\Theta^\dagger\Theta]$ is $SU(3)$.
Its Lie algebra is $\mathfrak{su}(3)$ with generators
$\{T^a\}_{a=1}^8$ satisfying $[T^a, T^b] = if^{abc}T^c$ where
$f^{abc}$ are the $SU(3)$ structure constants.
\end{theorem}

\begin{proof}
See \texttt{appendix\_G\_internal\_color\_symmetry.tex}, Theorem~G.A.
Summary: $V_c \cong \mathbb{C}^3$ and $\mathrm{Aut}(V_c, \langle\cdot,\cdot\rangle)
= U(3)$ in general; the unitarity plus $\det = 1$ (from the UBT volume form)
restricts to $SU(3)$.  The generators are identified with the $8 = \dim_\mathbb{R}(\mathcal{B}) - \dim U(1)$
independent traceless Hermitian transformations of $V_c$.
\end{proof}

\begin{proposition}[G.B: Fundamental Representation --- PROVED]
\label{prop:G_B}
The field $\Theta$ restricted to $V_c$ transforms in the fundamental
representation $\mathbf{3}$ of $SU(3)$.  The components $\Theta_I, \Theta_J, \Theta_K$
(projections onto $I, J, K$ directions) are the quark color triplet.
\end{proposition}

\begin{proof}
See \texttt{appendix\_G\_internal\_color\_symmetry.tex}, Proposition~G.B.
The $SU(3)$ action $\mathcal{U}\in SU(3)$ acts on $V_c$ by
$(\Theta_I, \Theta_J, \Theta_K)^T \mapsto \mathcal{U}(\Theta_I, \Theta_J, \Theta_K)^T$.
This is the defining (\textbf{3}) representation.
\end{proof}

\begin{theorem}[G.C: Adjoint Representation --- PROVED]
\label{thm:G_C}
The gluon field $A_\mu = A_\mu^a T^a$ transforms in the adjoint
representation $\mathbf{8}$ of $SU(3)$.  The field strength tensor
$F_{\mu\nu}^a = \partial_\mu A_\nu^a - \partial_\nu A_\mu^a + f^{abc}A_\mu^b A_\nu^c$
is reproduced from the internal phase curvature of $\Theta$.
\end{theorem}

\begin{proof}
See \texttt{appendix\_G\_internal\_color\_symmetry.tex}, Theorem~G.C.
The 8 generators of $\mathfrak{su}(3)$ act on $V_c$ as $3\times 3$
traceless Hermitian matrices (Gell-Mann matrices $\lambda^a$).  The
gauge field $A_\mu = A_\mu^a T^a$ is the connection on the principal
$SU(3)$-bundle over spacetime, with curvature $F_{\mu\nu}$.
\end{proof}

\begin{theorem}[G.D: Separation from Electroweak --- PROVED]
\label{thm:G_D}
The $SU(3)_c$ sector decouples from the electroweak $SU(2)_L\times U(1)_Y$
sector at the level of the UBT action $S[\Theta]$.
The direct product structure $SU(3)_c \times SU(2)_L \times U(1)_Y$ follows.
\end{theorem}

\begin{proof}
See \texttt{appendix\_G\_internal\_color\_symmetry.tex}, Theorem~G.D.
$SU(3)_c$ acts on $V_c$ (internal quaternionic directions $\{I,J,K\}$)
while $SU(2)_L$ acts on the external $\mathrm{Mat}(2,\mathbb{C})$ structure
and $U(1)_Y$ acts on the overall phase.  These actions commute because
the internal phase manifold $\mathcal{F}$ (compact, fiber over spacetime)
is independent of the spacetime fiber.
\end{proof}

%──────────────────────────────────────────────────────────────────────────────
\section{Comparison of Approaches}
\label{sec:comparison}
%──────────────────────────────────────────────────────────────────────────────

\begin{center}
\begin{tabular}{p{3.5cm}lp{5cm}}
  \toprule
  \textbf{Approach} & \textbf{Status} & \textbf{Notes} \\
  \midrule
  Three imaginary axes $\{I,J,K\}$ directly
    & Dead end
    & Gives $\mathfrak{su}(2)$, not $\mathfrak{su}(3)$ \\
  Octonion completion $\mathbb{O}$
    & Ruled out
    & $\mathbb{O}$ is non-associative; UBT requires associativity \\
  \textbf{Involution approach}
  $V_c = \ker(P_2+1)$
    & \textbf{Proved}
    & Gives $SU(3)$ from $8 = \dim_\mathbb{R}(\mathbb{C}\otimes\mathbb{H})$;
      Theorems G.A--G.D \\
  \bottomrule
\end{tabular}
\end{center}

%──────────────────────────────────────────────────────────────────────────────
\section{The Remaining Gap: Color Confinement}
\label{sec:confinement}
%──────────────────────────────────────────────────────────────────────────────

\begin{gapbox}
\textbf{Remaining gap: Color Confinement.}

Status: \textbf{Open} (Clay Millennium Prize Problem, unsolved since 2000).

UBT provides a \emph{qualitative} mechanism:
\begin{itemize}
  \item The non-Abelian curvature of the internal phase bundle $A_\mu$
        over spacetime creates an area law for Wilson loops:
        \[W(C) = \mathrm{Tr}\,\mathcal{P}\exp\!\Bigl(i\oint_C A_\mu dx^\mu\Bigr)
          \;\sim\; e^{-\sigma\,\mathrm{Area}(C)}\]
        for large loops $C$, where $\sigma$ is the string tension.
  \item This forbids isolated color charges (quarks) in the vacuum sector.
\end{itemize}

\textbf{Quantitative derivation} (from the UBT action $S[\Theta]$):
this is a non-perturbative computation requiring lattice-QCD-level
methods applied to the internal phase bundle.  No analytical proof
exists in the Standard Model; UBT is not expected to provide one
immediately.

\textbf{UBT contribution:} The geometric framework (internal phase bundle
with non-Abelian curvature) provides a natural language for confinement.
A quantitative proof would simultaneously solve the Clay Millennium Problem.
\end{gapbox}

%──────────────────────────────────────────────────────────────────────────────
\section{Derivation Status Summary}
\label{sec:status}
%──────────────────────────────────────────────────────────────────────────────

\begin{resultbox}
\textbf{SU(3) derivation status: PROVED [L0] with one open gap.}

\begin{center}
\begin{tabular}{lll}
  \toprule
  \textbf{Result} & \textbf{Status} & \textbf{Reference} \\
  \midrule
  Lie algebra $\mathfrak{su}(3)$ & \textbf{Proved} [L0] & Thm.~\ref{thm:G_A} \\
  Fundamental rep.\ (quarks, $\mathbf{3}$) & \textbf{Proved} [L0] & Prop.~\ref{prop:G_B} \\
  Adjoint rep.\ (gluons, $\mathbf{8}$) & \textbf{Proved} [L0] & Thm.~\ref{thm:G_C} \\
  Decoupling from EW & \textbf{Proved} [L0] & Thm.~\ref{thm:G_D} \\
  Color confinement & \textbf{Open} & Clay Prize; see §\ref{sec:confinement} \\
  \bottomrule
\end{tabular}
\end{center}

\textbf{DERIVATION\_INDEX.md update:}
Change ``$SU(3)_c$: Semi-empirical / Candidate construction'' to\\
``$SU(3)_c$: Proved [L0] --- Theorems G.A--G.D in \texttt{appendix\_G};
confinement gap remains (Millennium Problem).''
\end{resultbox}

\ifdefined\INCLUDEMODE
\else
  \end{document}
\fi
